% Part: proof-theory
% Chapter: normalization
% Section: introduction

\documentclass[../../../include/open-logic-section]{subfiles}

\begin{document}

\olfileid[pt]{pt}{nor}{int}

\olsection{Introdução}

Se você !!{prove}d uma condicional~$!A \lif !B$, talvez possa usá-la
em !!a{proof} de~$!B$ mediante~\Elim{\lif}. Isso, naturalmente, também
exigiria !!a{proof}~$\delta_1$ de~$!A$. É bastante provável que sua
!!{proof}~$\delta_1$ de~$!A \lif !B$ termine em~\Intro{\lif}, que
transforma !!a{proof}~$\delta_2$ de~$!B$ a partir da hipótese aberta~$!A$
em uma demonstração de~$!A \lif !B$. Nesse caso, sua !!{proof} final tem
a forma
\[
\AxiomC{$\Discharge{!A}{x}$}
\RightLabel{$\delta_2$}
\DeduceC{$!B$}
\DischargeRule{\Intro{\lif}}{x}
\UnaryInfC{$!A \lif !B$}
\AxiomC{}
\RightLabel{$\delta_2$}
\DeduceC{$!A$}
\RightLabel{\Elim{\lif}}
\BinaryInfC{$!B$}
\DisplayProof
\]
Embora sua estratégia seja eficiente por reutilizar algo que você já
tinha (!!a{proof} de~$!A \lif !B$), sua !!{proof} não é. Parece um
desvio introduzir $!A \lif !B$ apenas para eliminá-la imediatamente.
Deve haver uma maneira mais direta de !!{prove}~$!B$.

Na verdade, sempre há. ``Desvios'' como o acima---!!a{introduction}
imediatamente seguida por !!a{elimination}---sempre podem ser removidos
de !!{proof}s de dedução natural. Esse processo se chama
\emph{normalização}, e o resultado é uma !!{proof} \emph{normal} (ou
!!a{proof} \emph{em forma normal}).

A demonstração desse resultado, assim como a do teorema da eliminação do
corte no cálculo de sequentes, é um tanto trabalhosa e exige considerar
muitos casos. É mais difícil demonstrá-lo para o sistema clássico
\Log{N1c} do que para o sistema intuicionista \Log{N1i}. Assim como a
eliminação do corte mostra que não se pode !!{prove} mais usando a regra
\Cut{} no cálculo de sequentes, a normalização mostra que não se pode
demonstrar mais resultados usando desvios do que evitando-os. Há outras
semelhanças: uma !!{proof} normal pode ser muito mais longa do que a
!!{proof} mais curta, com desvios, do mesmo resultado, assim como
!!{proof}s no cálculo de sequentes podem ser muito mais curtas com cortes
do que a !!{proof} sem cortes mais curta. A propriedade da
sub-!!{formula}---segundo a qual sempre se pode !!{prove} um resultado
usando apenas sub-!!{formula}s desse resultado na !!{proof}---decorre da
normalização para a dedução natural assim como decorre da eliminação do
corte para o cálculo de sequentes.

\end{document}
